\section{Discussion}
\label{sec:discussion}

\subsection{Interpreting Performance--Energy Trade-offs}
\label{sec:discussion-tradeoffs}

The energy efficiency of an LLM serving architecture is determined jointly by average power and execution time. A deployment that lowers GPU utilization or operating frequency may draw less power but consume more total energy if communication, synchronization, or slower computation substantially extends the measurement interval. Conversely, a high-power configuration can be energy efficient when it completes requests quickly and returns accelerators to an idle or reusable state. Therefore, architecture comparisons should not rank systems by power, latency, or throughput in isolation. We interpret a configuration as preferable only when it satisfies the same service-level objective (SLO) and improves normalized energy metrics such as joules per successful request or tokens per joule.

This distinction is particularly important for disaggregated serving. PD and AF can specialize worker pools for phases or operators with different resource demands, but they also introduce data movement and coordination. The relevant question is not whether disaggregation reduces the power of an individual pool, but whether the reduction in computation or interference outweighs transfer energy, synchronization delay, pipeline bubbles, and the baseline power of additional active GPUs. Reporting per-role energy alongside deployment-level energy helps identify whether an apparent gain comes from genuine efficiency or from shifting work to another pool.

\subsection{Implications for Deployment Design}
\label{sec:discussion-implications}

No serving architecture should be expected to dominate across all workloads and platforms. Prefill-heavy workloads expose more compute parallelism, whereas decode-heavy and long-context workloads place greater pressure on model and KV-cache bandwidth. Native serving avoids inter-pool transfers but couples components with different scaling behavior. PD can independently provision compute-oriented prefill and memory-oriented decode, while AF can specialize attention and FFN resources, particularly for MoE models. The preferred architecture can consequently change with the input/output ratio, offered load, model type, and SLO.

Communication and parallelism must be considered as part of the architecture rather than as independent implementation details. TP and EP add frequent collectives or token exchange; PD transfers request state and KV data; and AF exchanges activations at layer boundaries. Fast scale-up links can hide or amortize some of these costs, whereas PCIe or multi-node RDMA may move the energy--performance crossover. Similarly, increasing DP can raise throughput through replication but also duplicates model memory and idle power. Practical provisioning should therefore select the architecture, role placement, and DP/TP/EP degrees jointly.

Hardware selection creates another trade-off. Datacenter GPUs provide large ECC memory and high-speed interconnects, while consumer GPUs can provide high arithmetic throughput at different capacity, management, and connectivity points. Peak FP16/BF16 throughput alone does not predict serving efficiency: decode can remain memory-bound, a model may not fit at the desired parallel degree, and a disaggregated design may depend more on communication than arithmetic capability. Frequency controls should likewise be applied selectively. Compute-bound components are more sensitive to core-frequency reductions, whereas memory- or communication-stalled components may tolerate them. Memory-frequency reductions are useful only when their power savings exceed the additional service time and do not violate the SLO.

\subsection{Measurement Scope and Threats to Validity}
\label{sec:discussion-validity}

Our conclusions are bounded by the evaluated models, request shapes, hardware, software version, and deployment configurations. The fixed-shape workloads isolate input/output effects and improve repeatability, but production traffic includes correlations, multi-turn context, shared prefixes, burstiness, cancellations, and changing output distributions. Open-loop Poisson arrivals do not reproduce every production process. Results should therefore be interpreted as controlled architectural characterization rather than a complete model of every deployment.

Hardware and software effects also limit generalization. GPU firmware, drivers, kernels, quantization, attention backends, routing policies, and batching algorithms can change utilization and energy. Product-level peak specifications do not guarantee realized throughput. We mitigate these effects by keeping the software and precision fixed within paired comparisons, recording manifests and system state, validating communication paths, and repeating deployments; nevertheless, results from the evaluated GPUs should not be directly extrapolated to accelerators with different memory hierarchies or interconnects.

The energy boundary is another limitation. Our primary metric includes the GPUs assigned to a deployment because their cumulative energy counters provide a consistent measurement across configurations. It does not include host CPU, DRAM, NIC, switch, storage, cooling, or power-conversion losses. This boundary is suitable for comparing GPU execution but can understate the cost of multi-node and communication-intensive architectures. Future work should combine device counters with node- and rack-level metering and should account for embodied carbon and regional electricity carbon intensity.

Finally, SLO filtering can introduce survivorship bias if only successful high-load configurations remain. We therefore report offered and achieved load, success rate, and latency distributions and compare energy only among configurations satisfying the same validity criteria. Repeated trials quantify run-to-run variation, but longer-duration studies are still needed to capture thermal drift, failures, autoscaling transients, and diurnal workloads.

\subsection{Toward Architecture-Aware Energy Management}
\label{sec:discussion-future}

The measurement dimensions in this study suggest that energy management should be architecture aware. Rather than tuning frequency, batch size, parallelism, or placement independently, a controller could use workload composition and measured phase/operator sensitivity to choose a complete operating point. Such a controller would select among Native, PD, and AF; assign hardware to each role; determine DP/TP/EP degrees; and apply role-specific core or memory settings subject to TTFT, TPOT, capacity, and power constraints. The resulting objective is a constrained Pareto optimization problem rather than simple power minimization.

A useful next step is to construct predictive models from the collected measurements. These models should capture computation, memory traffic, communication, synchronization, and idle power separately and estimate both latency and energy before deployment. Online feedback could then correct prediction error and adapt to workload drift. Extending the methodology to multi-tenant serving, heterogeneous accelerators, shared KV-cache tiers, node-level energy measurement, and carbon-aware placement would provide a more complete foundation for sustainable LLM infrastructure.
